% GENERATED FILE. Edit canonical lessons, then run npm run generate:books.
% canonical-source-hash: fnv1a64:c83a77713117ceb9
% canonical-lessons: SA-C452-udvegah, SA-C452-hatasah, SA-C452-visidati, SA-C452-vyagrah, SA-C452-dinah

\chapter{Worry and Despair}
\label{ch:sa-worry-and-despair}
\hlchaptermodality{452}

\begin{chapteropening}
\textbf{By the end of this chapter:} \emph{I can say that someone is worried or has lost hope.}

Complete the last lesson of chapter 452: I can say that someone is worried or has lost hope.
\end{chapteropening}

\section[\texorpdfstring{\sk{उद्वेगः} (udvegaḥ)}{उद्वेगः (udvegaḥ)}]{\sk{उद्वेगः} (udvegaḥ) — anxiety}

\label{lesson:SA-C452-udvegah}



\begin{quote}
Before the new one: say the Sanskrit for troubled, then the Sanskrit for anguish.
\end{quote}

\subsection*{You'll want to know: \sk{उद्वेगः}}
\textbf{\sk{उद्वेगः}} — \emph{udvegaḥ} — \textquotedblleft{}anxiety\textquotedblright{}.

\begin{quote}
\sk{परीक्षायाः} \sk{पूर्वं} \sk{छात्रस्य} \sk{उद्वेगः} \sk{भवति।}
\end{quote}

\begin{quote}
\emph{parīkṣāyāḥ pūrvaṃ chātrasya udvegaḥ bhavati.}
\end{quote}

\begin{quote}
A student feels anxious before an exam.
\end{quote}

A masculine noun.

\begin{cousinweb}
\textbf{\sk{उद्वेगः}} (\emph{udvegaḥ}) is built from \emph{ud}, \textquotedblleft{}up\textquotedblright{}, and the root √vij, \textquotedblleft{}to tremble\textquotedblright{}; MW gives \textquotedblleft{}trembling\textquotedblright{} and \textquotedblleft{}agitation, anxiety\textquotedblright{} (MW).
\end{cousinweb}

\begin{grammarlens}[title={What you've built}]
One more word for describing what you lived and felt.
\end{grammarlens}

\subsection*{Your turn}
\begin{itemize}
\raggedright
  \item \emph{Say it:} \emph{udvegaḥ}
  \item \emph{Say it:} \emph{udvegaḥ}, once more
  \item \emph{Say it:} \emph{parīkṣāyāḥ pūrvaṃ chātrasya udvegaḥ bhavati.} — A student feels anxious before an exam.
  \item \emph{Recall:} say the Sanskrit for troubled, then the Sanskrit for anguish, then say \emph{udvegaḥ} again
\end{itemize}

\begin{tcolorbox}[breakable,colback=teal!4,colframe=teal!35!black,title={Before you move on}]
What does \sk{उद्वेगः} mean? (\textquotedblleft{}Anxiety\textquotedblright{}.) Say it once more.
\end{tcolorbox}
\section[\texorpdfstring{\sk{हताशः} (hatāśaḥ)}{हताशः (hatāśaḥ)}]{\sk{हताशः} (hatāśaḥ) — despairing; without hope}

\label{lesson:SA-C452-hatasah}



\begin{quote}
Before the new one: say the Sanskrit for anguish, then the Sanskrit for anxiety.
\end{quote}

\subsection*{You'll want to know: \sk{हताशः}}
\textbf{\sk{हताशः}} — \emph{hatāśaḥ} — \textquotedblleft{}despairing; without hope\textquotedblright{}.

\begin{quote}
\sk{वर्षा} \sk{न} \sk{भवति}, \sk{अतः} \sk{कृषकः} \sk{हताशः} \sk{अस्ति।}
\end{quote}

\begin{quote}
\emph{varṣā na bhavati, ataḥ kṛṣakaḥ hatāśaḥ asti.}
\end{quote}

\begin{quote}
There is no rain, so the farmer has lost hope.
\end{quote}

\begin{cousinweb}
\textbf{\sk{हताशः}} (\emph{hatāśaḥ}) is built from \emph{hata}, \textquotedblleft{}struck, destroyed\textquotedblright{}, from the root √han, and \textbf{\sk{आशा}}, \textquotedblleft{}hope\textquotedblright{}: one whose hopes are struck down; MW gives \textquotedblleft{}whose hopes are destroyed, desperate\textquotedblright{} (MW).
\end{cousinweb}

\begin{grammarlens}[title={What you've built}]
One more word for describing what you lived and felt.
\end{grammarlens}

\subsection*{Your turn}
\begin{itemize}
\raggedright
  \item \emph{Say it:} \emph{hatāśaḥ}
  \item \emph{Say it:} \emph{hatāśaḥ}, once more
  \item \emph{Say it:} \emph{varṣā na bhavati, ataḥ kṛṣakaḥ hatāśaḥ asti.} — There is no rain, so the farmer has lost hope.
  \item \emph{Recall:} say the Sanskrit for anguish, then the Sanskrit for anxiety, then say \emph{hatāśaḥ} again
\end{itemize}

\begin{tcolorbox}[breakable,colback=teal!4,colframe=teal!35!black,title={Before you move on}]
What does \sk{हताशः} mean? (\textquotedblleft{}Despairing; without hope\textquotedblright{}.) Say it once more.
\end{tcolorbox}
\section[\texorpdfstring{\sk{विषीदति} (viṣīdati)}{विषीदति (viṣīdati)}]{\sk{विषीदति} (viṣīdati) — loses heart}

\label{lesson:SA-C452-visidati}



\begin{quote}
Before the new one: say the Sanskrit for anxiety, then the Sanskrit for despairing.
\end{quote}

\subsection*{You'll want to know: \sk{विषीदति}}
\textbf{\sk{विषीदति}} — \emph{viṣīdati} — \textquotedblleft{}loses heart\textquotedblright{}.

\begin{quote}
\sk{कार्यं} \sk{कठिनम्} \sk{अस्ति}, \sk{किन्तु} \sk{अरुणः} \sk{न} \sk{विषीदति।}
\end{quote}

\begin{quote}
\emph{kāryaṃ kaṭhinam asti, kintu aruṇaḥ na viṣīdati.}
\end{quote}

\begin{quote}
The work is hard, but Arun does not lose heart.
\end{quote}

Past (he …): \sk{सः} \sk{विषण्णवान्}, \emph{saḥ viṣaṇṇavān}.

\begin{cousinweb}
\textbf{\sk{विषीदति}} (\emph{viṣīdati}) is built from \emph{vi} and the root √sad, \textquotedblleft{}to sit, to sink\textquotedblright{}, the verb behind \textbf{\sk{विषादः}}; MW gives \textquotedblleft{}to be dejected, despond, despair\textquotedblright{} (MW).
\end{cousinweb}

\begin{grammarlens}[title={What you've built}]
One more word for describing what you lived and felt.
\end{grammarlens}

\subsection*{Your turn}
\begin{itemize}
\raggedright
  \item \emph{Say it:} \emph{viṣīdati}
  \item \emph{Say it:} \emph{viṣīdati}, once more
  \item \emph{Say it:} \emph{kāryaṃ kaṭhinam asti, kintu aruṇaḥ na viṣīdati.} — The work is hard, but Arun does not lose heart.
  \item \emph{Recall:} say the Sanskrit for anxiety, then the Sanskrit for despairing, then say \emph{viṣīdati} again
\end{itemize}

\begin{tcolorbox}[breakable,colback=teal!4,colframe=teal!35!black,title={Before you move on}]
What does \sk{विषीदति} mean? (\textquotedblleft{}Loses heart\textquotedblright{}.) Say it once more.
\end{tcolorbox}
\section[\texorpdfstring{\sk{व्यग्रः} (vyagraḥ)}{व्यग्रः (vyagraḥ)}]{\sk{व्यग्रः} (vyagraḥ) — agitated; distracted}

\label{lesson:SA-C452-vyagrah}



\begin{quote}
Before the new one: say the Sanskrit for despairing, then the Sanskrit for loses heart.
\end{quote}

\subsection*{You'll want to know: \sk{व्यग्रः}}
\textbf{\sk{व्यग्रः}} — \emph{vyagraḥ} — \textquotedblleft{}agitated; distracted\textquotedblright{}.

\begin{quote}
\sk{विलम्बेन} \sk{यात्री} \sk{व्यग्रः} \sk{भवति।}
\end{quote}

\begin{quote}
\emph{vilambena yātrī vyagraḥ bhavati.}
\end{quote}

\begin{quote}
The passenger grows agitated at the delay.
\end{quote}

\begin{cousinweb}
\textbf{\sk{व्यग्रः}} (\emph{vyagraḥ}) is built from \emph{vi}, \textquotedblleft{}apart\textquotedblright{}, and \emph{agra}, \textquotedblleft{}a point\textquotedblright{}: a mind with no single point; MW gives \textquotedblleft{}distracted\textquotedblright{} and \textquotedblleft{}bewildered, agitated, alarmed\textquotedblright{} (MW).
\end{cousinweb}

\begin{grammarlens}[title={What you've built}]
One more word for describing what you lived and felt.
\end{grammarlens}

\subsection*{Your turn}
\begin{itemize}
\raggedright
  \item \emph{Say it:} \emph{vyagraḥ}
  \item \emph{Say it:} \emph{vyagraḥ}, once more
  \item \emph{Say it:} \emph{vilambena yātrī vyagraḥ bhavati.} — The passenger grows agitated at the delay.
  \item \emph{Recall:} say the Sanskrit for despairing, then the Sanskrit for loses heart, then say \emph{vyagraḥ} again
\end{itemize}

\begin{tcolorbox}[breakable,colback=teal!4,colframe=teal!35!black,title={Before you move on}]
What does \sk{व्यग्रः} mean? (\textquotedblleft{}Agitated; distracted\textquotedblright{}.) Say it once more.
\end{tcolorbox}
\section[\texorpdfstring{\sk{दीनः} (dīnaḥ)}{दीनः (dīnaḥ)}]{\sk{दीनः} (dīnaḥ) — wretched; dejected}

\label{lesson:SA-C452-dinah}



\begin{quote}
Before the new one: say the Sanskrit for loses heart, then the Sanskrit for agitated.
\end{quote}

\subsection*{You'll want to know: \sk{दीनः}}
\textbf{\sk{दीनः}} — \emph{dīnaḥ} — \textquotedblleft{}wretched; dejected\textquotedblright{}.

\begin{quote}
\sk{दीनः} \sk{याचकः} \sk{मार्गे} \sk{तिष्ठति।}
\end{quote}

\begin{quote}
\emph{dīnaḥ yācakaḥ mārge tiṣṭhati.}
\end{quote}

\begin{quote}
A wretched beggar stands by the road.
\end{quote}

\begin{cousinweb}
\textbf{\sk{दीनः}} (\emph{dīnaḥ}) MW asks, with a question mark, whether it comes from a root √dī, so its origin is uncertain; its first sense is \textquotedblleft{}scarce, scanty\textquotedblright{}, then \textquotedblleft{}depressed, afflicted, sad\textquotedblright{} and \textquotedblleft{}miserable, wretched\textquotedblright{} (MW).
\end{cousinweb}

\begin{grammarlens}[title={What you've built}]
One more word for describing what you lived and felt.
\end{grammarlens}

\subsection*{Your turn}
\begin{itemize}
\raggedright
  \item \emph{Say it:} \emph{dīnaḥ}
  \item \emph{Say it:} \emph{dīnaḥ}, once more
  \item \emph{Say it:} \emph{dīnaḥ yācakaḥ mārge tiṣṭhati.} — A wretched beggar stands by the road.
  \item \emph{Recall:} say the Sanskrit for loses heart, then the Sanskrit for agitated, then say \emph{dīnaḥ} again
\end{itemize}

\begin{tcolorbox}[breakable,colback=teal!4,colframe=teal!35!black,title={Before you move on}]
What does \sk{दीनः} mean? (\textquotedblleft{}Wretched; dejected\textquotedblright{}.) Say it once more.
\end{tcolorbox}
